\section{Corpus Construction}
\label{sec:dataset}

\subsection{Source and document structure}

The corpus is drawn from the U.S. Securities and Exchange Commission's EDGAR
system \citep{sec2024edgar}, which publishes every mandatory disclosure filed by
every U.S. public company. We ingested twelve monthly XBRL feeds labelled
2025-09 through 2026-08. Filing dates in the resulting corpus span 2025-08 to
2026-08, because a monthly feed includes filings dated at the end of the preceding
month; \Cref{fig:months} shows the distribution, with the small 2025-08 bucket
being that spillover.

The structural point that motivates this paper is that an annual report is not one
document. A Form 10-K is a container: a mandated sequence of numbered
\emph{Items}, plus attached \emph{Exhibits}, filed under a single accession
number. Three containers matter here.

\textbf{Item 1A, Risk Factors}, is narrative prose in which a company enumerates
what could go wrong. It is the longest narrative section, averaging 116{,}779
characters and reaching 1{,}060{,}461 in this corpus (\Cref{tab:items}). It is the
only section we embed.

\textbf{Exhibit 21} is an attachment, not a section, required by Item 601(b)(21)
of Regulation S-K. It lists subsidiaries and their jurisdiction of incorporation,
so it parses into rows rather than prose. Two properties matter. First, the
distribution is extremely skewed: 275 filers disclose a single subsidiary while
62 disclose more than 500, and the largest discloses 2{,}780
(\Cref{fig:subdist}). Second, and more consequentially,
\emph{EX-21 is incomplete by design}: rule 601(b)(21)(ii) permits a filer to omit
any subsidiary that is not a ``significant subsidiary''. The exhibit is therefore
a floor on corporate structure and never a census, and the graph inherits that
limit. An absent subsidiary is not evidence that none exists.

\textbf{Forms 3, 4 and 5} are separate filings, not part of the 10-K. Officers,
directors and ten-percent owners report their positions and transactions in
structured XML. This is the container that makes identity resolvable, because each
filing carries the insider's own Central Index Key (CIK) --- a stable SEC
identifier. Two directors who share a surname are separable on CIK even though
their names are not.

\Cref{fig:containers} summarizes the arrangement and the measured disjointness
that follows from it.

\subsection{Ingestion}

The pipeline runs in nine stages: fetch and parse the monthly feeds; fetch the
documents each filing manifests; extract narrative Items from HTML; parse EX-21
and the ownership forms; stage to object storage; load the relational store;
embed; and build indexes and the graph projection. \Cref{fig:pipeline} gives the
wall-clock breakdown, totalling roughly 252 minutes. Two stages dominate:
document fetch (118 minutes) and embedding (87 minutes), together about 82\% of
the total. The fetch cost is externally bounded --- EDGAR's fair-access policy
caps automated requests at ten per second per requester
\citep{sec2024edgar} --- so no amount of parallelism removes it.

The filings manifest 1{,}904{,}128 documents in total. We fetched and parsed the
10{,}592 that carry the narrative and exhibit content the benchmark needs,
amounting to 19.7\,GB of raw HTML, which reduces to a 3{,}803\,MiB relational
store plus a 1.0\,GB vector index. \Cref{tab:corpus} gives the resulting counts.

\subsection{Schema}

The relational store is PostgreSQL. Filings and companies are keyed by accession
number and CIK. \texttt{filing\_section} holds one row per extracted Item with its
text; \texttt{section\_chunk} holds the embedded chunks; \texttt{subsidiary} holds
parsed EX-21 rows; \texttt{reporting\_owner} holds parsed ownership rows; and
\texttt{filing\_auditor} holds auditor attestations.

The graph is a projection of the same data, not a second source of truth. It holds
430{,}257 nodes across five labels --- \texttt{Filing} (167{,}050),
\texttt{Subsidiary} (200{,}869), \texttt{Person} (51{,}558), \texttt{Company}
(10{,}511) and \texttt{AuditFirm} (269) --- and 561{,}144 edges across seven
types: \texttt{:SUBSIDIARY\_OF} (212{,}124), \texttt{:FILED} (167{,}050),
\texttt{:OFFICER\_OF} (86{,}180), \texttt{:DIRECTOR\_OF} (70{,}167),
\texttt{:OWNS\_SHARES} (19{,}542), \texttt{:AUDITED\_BY} (5{,}835) and
\texttt{:RESOLVES\_TO} (246).

Two design decisions in that schema carry most of the paper's weight.

\emph{Every edge carries the accession number of the filing that asserts it.}
Provenance is a property of the relationship rather than of a document, which is
what lets an answer cite the exact filing supporting each individual claim rather
than the set of documents that were retrieved.

\emph{\texttt{:RESOLVES\_TO} is the entity-resolution layer, and it is thin.}
It links a subsidiary named in an EX-21 to a company that files in its own right,
a fact no single document states. There are 246 such edges against 200{,}869
subsidiary nodes --- a coverage of 0.12\% --- because EX-21 carries no CIK for a
subsidiary, so the link must be established by name matching. We show in
\Cref{sec:failures} that this thinness is a direct cause of the graph losing the
one category built to exercise it.
